\section{Zero transport, real births, and complex event residues}
\label{tr:section}

This section distinguishes three quantities that agree in some double-zero
models but are different in general: the number of complex zeros counted
with multiplicity, the number of real zeros on a regular time slice, and
the residue charge of a multiple-zero event.  Every domain below is finite.
In particular, no limiting contour, infinite product over zeros, or
unproved assertion about zeros at arbitrarily large height is used.
We use the backward-heat normalization
\begin{equation}\label{tr:heat}
                    H_t=-H_{zz}.
\end{equation}
When real variables are used, $H(t,x)$ is real and $H_x=H_z$ on the real
slice.  Joint real analyticity permits local complexification.  All
arguments involving complex time explicitly assume joint holomorphy on
a neighborhood of the indicated closed product.

\subsection{The local heat polynomial and its roots}

\begin{lemma}[Elementary Hermite root calculation]\label{tr:hermite-roots}
For $m\geq1$ define
\[
 p_m(s,w)=\sum_{k=0}^{\lfloor m/2\rfloor}
        \frac{m!(-s)^k w^{m-2k}}{k!(m-2k)!}.
\]
Then $p_m(1,\cdot)$ has $m$ distinct real roots.  The polynomial
$p_m(-1,\cdot)$ has no real root when $m$ is even and exactly one real
root, a simple root at zero, when $m$ is odd.  If
$\lambda_1,\ldots,\lambda_m$ are the roots of the monic probabilists'
Hermite polynomial, then the roots of $p_m(1,\cdot)$ are
$\sqrt2\lambda_1,\ldots,\sqrt2\lambda_m$.
\end{lemma}

\begin{proof}
Define
\[
 \operatorname{He}_m(u)=(-1)^m e^{u^2/2}
                       \frac{d^m}{du^m}e^{-u^2/2}.
\]
Differentiating $e^{u\xi-\xi^2/2}$ in $\xi$, or expanding its two
exponentials, gives
\[
 \operatorname{He}_m(u)=
 \sum_{k=0}^{\lfloor m/2\rfloor}
 \frac{m!(-1)^k u^{m-2k}}{2^k k!(m-2k)!},
 \qquad p_m(1,w)=2^{m/2}\operatorname{He}_m(w/\sqrt2).
\]
Repeated integration by parts, whose boundary terms vanish by Gaussian
decay, shows that
\[
 \int_{\mathbb R}\operatorname{He}_m(u)q(u)e^{-u^2/2}\,du=0
                      \quad(\deg q<m).
\]
Suppose that $\operatorname{He}_m$ has fewer than $m$ distinct real zeros
of odd multiplicity.  Let $q$ be the product of the linear factors at
those zeros, once each.  Then $\deg q<m$, while
$\operatorname{He}_m q$ has constant sign on the real line apart from
its zeros and is not identically zero.  Its Gaussian integral cannot
vanish, a contradiction.  Degree $m$ forces all its roots to be real and
simple.  Finally $p_m(-1,w)$ is a polynomial in $w^2$ with strictly
positive coefficients if $m$ is even; if $m$ is odd it is $w$ times such
a polynomial, with a strictly positive constant term.  This proves the
negative-time assertion as well.
\end{proof}

\begin{theorem}[Local analytic splitting and isolation]\label{tr:splitting}
Suppose $H$ is holomorphic near $(t_*,z_*)$, satisfies
\eqref{tr:heat}, and
\[
 H(t_*,z_*+w)=c_m w^m+O(w^{m+1}),\qquad c_m\ne0.
\]
There is a proper local cluster of exactly $m$ spatial zeros, counted
with multiplicity.  With $s=t-t_*\ne0$ sufficiently small, all its roots
are simple.  More precisely, after $s=\varepsilon^2$, they can be written
\begin{equation}\label{tr:root-expansion}
 z_j(t_*+\varepsilon^2)-z_*
       =\varepsilon\bigl(\sqrt2\lambda_j+O(\varepsilon)\bigr),
       \qquad 1\leq j\leq m,
\end{equation}
with the parenthesized functions holomorphic in $\varepsilon$.
For $m\geq2$, $(t_*,z_*)$ is consequently an isolated common zero of
$(H,H_z)$ in complex two-space.

If $t_*,z_*$ are real and $H$ is of real type, all $m$ roots are real
for small positive $s$.  For small negative $s$, exactly $m\bmod2$ of
them are real.  Thus the local number of real zeros increases by
$2\lfloor m/2\rfloor$ as real time passes through the event.
\end{theorem}

\begin{proof}
Write $H(t_*,z_*+w)=\sum_{n\geq m}c_nw^n$.  The heat equation identifies
the convergent bivariate Taylor coefficients as
\[
 H(t_*+s,z_*+w)=\sum_{k,l\geq0}
       \frac{(-1)^l(k+2l)!}{k!l!}c_{k+2l}s^lw^k.
\]
Every nonzero term has weighted degree at least $m$, when $s$ has weight
two and $w$ weight one.  Substitution $s=\varepsilon^2$, $w=\varepsilon y$
therefore gives a holomorphic extension to $\varepsilon=0$ of
\[
 F(\varepsilon,y)=\varepsilon^{-m}
                  H(t_*+\varepsilon^2,z_*+\varepsilon y),
 \qquad F(0,y)=c_m p_m(1,y).
\]
Convergence is uniform on every sufficiently small $\varepsilon$-disk
and every fixed bounded $y$-disk.  Lemma~\ref{tr:hermite-roots} and the
holomorphic implicit function theorem give $m$ distinct root branches
$y_j(\varepsilon)=\sqrt2\lambda_j+O(\varepsilon)$.

Choose a fixed small spatial circle around $z_*$ on which
$H(t_*,\cdot)$ is nonzero and whose interior contains just its order-$m$
zero.  Uniform continuity on that circle, followed by Rouch\'e's theorem,
preserves the count $m$ for all sufficiently small complex $s$.
The displayed branches lie inside this circle and account for every
root.  Since $F_y$ remains nonzero at them, $H_z$ is nonzero whenever
$\varepsilon\ne0$.  This proves isolation in complex two-space.

For real positive $\varepsilon$, each branch is real by uniqueness and
complex conjugation, so the positive-time conclusion follows.  At
negative real $s$, take $\varepsilon=i\sqrt{|s|}$.  If $\lambda_j\ne0$,
\eqref{tr:root-expansion} has a nonzero imaginary leading term of order
$\sqrt{|s|}$, and its error is $O(|s|)$, so that root is not real.
There is no zero Hermite root when $m$ is even.  When $m$ is odd there
is exactly one, because parity gives a zero and all Hermite roots are
simple.  For its branch, $w_0(\varepsilon)=\varepsilon y_0(\varepsilon)
=O(\varepsilon^2)$.  Both $w_0(\varepsilon)$ and $w_0(-\varepsilon)$ are
roots of the same function $H(t_*+\varepsilon^2,z_*+w)$ and correspond
to the unique $y$-root near zero: the second has normalized coordinate
$w_0(-\varepsilon)/\varepsilon=-y_0(-\varepsilon)=O(\varepsilon)$.
Hence they coincide.  The function
$w_0$ is even in $\varepsilon$ and is holomorphic in $s$.  Reality on
positive real $s$ makes its Taylor coefficients real, so it remains
real for negative real $s$.  It is the sole real root there.
\end{proof}

\subsection{An angular-current formula with every wall term}

Give the real $(t,x)$-plane the orientation $dt\wedge dx$.  Wherever
$J=(H,H_x)\ne0$, put
\begin{equation}\label{tr:angular-form}
 \alpha=\frac{H\,dH_x-H_x\,dH}{H^2+H_x^2},\qquad
 \beta=-\frac{\alpha}{2\pi}.
\end{equation}
Locally $\alpha=d\arg(H+iH_x)$, so $d\alpha=d\beta=0$ away from joint
zeros.  The sign in $\beta$ is chosen to assign a positive real charge
to a backward-heat event.  On a spatial cap its coefficient is
\begin{equation}\label{tr:cap-density}
 \beta\big|_{t=\mathrm{constant}}
      =\frac{H_x^2-HH_{xx}}{2\pi(H^2+H_x^2)}\,dx.
\end{equation}

\begin{lemma}[Phase increments and signed crossings]\label{tr:phase-jumps}
Let an oriented analytic path carry a nonzero real vector $(A,B)$, with
$A$ nonzero at its endpoints and not identically zero.  At an isolated
zero $v_j$ of $A$, put
\[
 q_j=\operatorname{sgn}B(v_j)\,
 \frac{\operatorname{sgn}A(v_j+\delta)
          -\operatorname{sgn}A(v_j-\delta)}2,
\]
where $\delta>0$ is small enough to isolate that zero.  Then
\begin{equation}\label{tr:path-phase}
 \int\frac{A\,dB-B\,dA}{A^2+B^2}
   =\left[\arctan(B/A)\right]_{\rm start}^{\rm end}
                       -\pi\sum_jq_j.
\end{equation}
Even-order tangencies contribute zero.  At a transverse crossing,
$q_j=\operatorname{sgn}(B A')(v_j)$.
\end{lemma}

\begin{proof}
Off the zeros of $A$, the differential of the principal arctangent
$\arctan(B/A)\in(-\pi/2,\pi/2)$ is the integrand.  At a zero of $A$,
$B$ is nonzero.  The one-sided limits of the arctangent have difference
$\pi q_j$: they are equal at an even tangency and differ by the stated
signed multiple of $\pi$ at an odd crossing.  Summing the integrals on
the complementary path intervals gives \eqref{tr:path-phase}.  The
formula is also the increment of a continuous argument of $A+iB$.
\end{proof}

\begin{proposition}[Cap count and wall correction]\label{tr:boundary-count}
Let
\[
 \mathcal R=\{(t,x):a<t<b,\quad \ell(t)<x<r(t)\}
\]
be a bounded moving rectangle, with analytic walls $\ell<r$ on a
neighborhood of $[a,b]$.  Assume $H$ is real analytic near its closure,
$J=(H,H_x)$ is nonzero on its whole boundary, and $H$ is nonzero at the
four corners.  Let $N_a,N_b$ count the zeros of $H$ on the two open
spatial caps; these zeros are simple.  For either wall $c=\ell,r$, set
$F_c(t)=H(t,c(t))$ and
\begin{equation}\label{tr:wall-flux}
 Q_c=\sum_{F_c(t_j)=0}\operatorname{sgn}H_x(t_j,c(t_j))\,
 \frac{\operatorname{sgn}F_c(t_j+\delta_j)
       -\operatorname{sgn}F_c(t_j-\delta_j)}2.
\end{equation}
The sums are finite, and
\begin{equation}\label{tr:count-winding}
 \frac1{2\pi}\oint_{\partial\mathcal R}\alpha
       =\frac{N_a-N_b-Q_\ell+Q_r}{2},\qquad
 \oint_{\partial\mathcal R}\beta
       =\frac{N_b-N_a+Q_\ell-Q_r}{2}.
\end{equation}
The orientation traverses the lower wall with increasing $t$, the late
cap with increasing $x$, the upper wall with decreasing $t$, and the
early cap with decreasing $x$.
\end{proposition}

\begin{proof}
Analyticity and nonzero corner values make each wall function nonzero
identically and give finitely many wall zeros on the compact interval.
At such a zero, boundary regularity gives $H_x\ne0$.  The same argument
on each cap gives finitely many simple spatial zeros.  Apply
Lemma~\ref{tr:phase-jumps} to each of the four edges with
$(A,B)=(H,H_x)$.  Along a forward cap, the signed crossing number is
$+1$ at every root, since $A'=H_x=B$.  Along a forward wall it is the
corresponding summand of $Q_c$.  The four endpoint arctangent
differences telescope.  The remaining contributions, with the declared
orientations, are $-\pi Q_\ell-\pi N_b+\pi Q_r+\pi N_a$.
This proves both identities.
\end{proof}

\begin{remark}[A single cap still counts ordinary zeros]\label{tr:cap-roots}
For fixed walls $\ell<r$ with $H(t,\ell)H(t,r)\ne0$, let
\[
 B(t)=\int_\ell^r\beta,
 \qquad\theta_c(t)=\arctan\frac{H_x(t,c)}{H(t,c)}.
\]
On a regular cap, Lemma~\ref{tr:phase-jumps} gives exactly
\begin{equation}\label{tr:single-cap}
 B(t)=\frac{N_t}{2}
                +\frac{\theta_\ell(t)-\theta_r(t)}{2\pi}.
\end{equation}
Thus simplicity of every cap root does not imply that the cap integral
vanishes.  If $H$ never vanishes on either wall, the wall contribution
to $\oint\beta$ is
$[\theta_r(b)-\theta_r(a)-\theta_\ell(b)+\theta_\ell(a)]/(2\pi)$;
it cancels precisely the endpoint terms of the two cap formulas.
\end{remark}

\begin{theorem}[Complete real birth balance]\label{tr:real-balance}
Assume the hypotheses of Proposition~\ref{tr:boundary-count}, the heat
equation \eqref{tr:heat}, and no identically zero spatial slice in the
analytic neighborhood.  The interior event set $\{H=H_x=0\}$ is finite.
If $m(T,x_*)$ is the exact spatial multiplicity at an event, then
\begin{align}
 \operatorname{ind}_{(T,x_*)}(H,H_x)
       &=-\left\lfloor\frac{m(T,x_*)}{2}\right\rfloor,
                                                        \label{tr:local-index}\\
 \oint_{\partial\mathcal R}\beta
       &=\sum_{(T,x_*)\in\mathcal R:\,H=H_x=0}
                   \left\lfloor\frac{m(T,x_*)}{2}\right\rfloor,
                                                        \label{tr:real-charge}\\
 N_b-N_a+Q_\ell-Q_r
       &=2\sum_{(T,x_*)\in\mathcal R:\,H=H_x=0}
                   \left\lfloor\frac{m(T,x_*)}{2}\right\rfloor.
                                                        \label{tr:birth-balance}
\end{align}
Here the local index is the counterclockwise winding of
$(H+iH_x)/|H+iH_x|$ on a small positively oriented loop.
\end{theorem}

\begin{proof}
Each event has finite spatial multiplicity by the nonzero-slice
hypothesis and is isolated by Theorem~\ref{tr:splitting}.  Boundary
regularity, compactness, and analyticity then make the event set finite:
an infinite set would have an accumulation point in the compact closure,
which could be neither a regular point nor an isolated event.

For one event, choose a small fixed real spatial interval enclosing its
local cluster and a small time interval symmetric about $T$.  Its walls
have $H\ne0$, its time caps are regular, and its sole event is the given
one.  Theorem~\ref{tr:splitting} gives cap counts $m\bmod2$ and $m$.
Proposition~\ref{tr:boundary-count} therefore gives local winding
$((m\bmod2)-m)/2=-\lfloor m/2\rfloor$, proving
\eqref{tr:local-index}.  Remove disjoint small neighborhoods of the
finitely many events from $\mathcal R$.  On the remainder $d\beta=0$.
Stokes' theorem, with the inner boundaries oppositely oriented, gives
\eqref{tr:real-charge}.  Combining it with
\eqref{tr:count-winding} proves \eqref{tr:birth-balance}.
\end{proof}

For a transverse wall crossing the formula can also be checked by
ordinary root motion.  A simple root satisfies
$x'(t)=H_{xx}/H_x$, so
\[
 Q_c=\sum_{H(t,c(t))=0}
       \operatorname{sgn}\{H_x(-H_{xx}+c'H_x)\}.
\]
A positive lower-wall crossing removes a root from the interval; a
positive upper-wall crossing adds one.  This is the sign convention in
\eqref{tr:birth-balance}.  The sign-change definition
\eqref{tr:wall-flux}, unlike this transverse formula, also handles
higher-order contacts.

\subsection{Conservation of the ordinary complex divisor}

\begin{proposition}[Multiplicity conservation on a proper spatial domain]
\label{tr:complex-conservation}
Let $\Omega$ be a bounded domain with piecewise smooth boundary.  Suppose
$H(t,\cdot)$ is holomorphic on a neighborhood of $\overline\Omega$,
continuously with its first spatial derivative in $t\in[a,b]$, and
$H(t,z)\ne0$ for $z\in\partial\Omega$.  Then
\begin{equation}\label{tr:complex-count}
 N_{\mathbb C}(t)=\frac1{2\pi i}\oint_{\partial\Omega}
                                      \frac{H_z(t,z)}{H(t,z)}\,dz
\end{equation}
is independent of $t$, with every root counted with its multiplicity.
No heat equation is needed.
\end{proposition}

\begin{proof}
The argument principle makes \eqref{tr:complex-count} an integer.  The
nonzero boundary condition and compactness give a positive common
boundary margin.  Hence the displayed integral is continuous in $t$,
and a continuous integer-valued function on an interval is constant.
\end{proof}

In particular, collision and splitting are not creation or destruction
of the complex divisor.  The $m$ descendants in
Theorem~\ref{tr:splitting} have the same total multiplicity as the
initial order-$m$ zero.  Their realness can change, which is why
\eqref{tr:birth-balance} carries an event term while
\eqref{tr:complex-count} does not.  A domain that permits zeros to cross
its spatial boundary must retain the signed boundary crossings; the
properness hypothesis is exactly what eliminates them here.

\subsection{Preparation and the discriminant residue}

Let $D\subset\mathbb C$ be a bounded disk or rectangle and
$\Omega_R=\{z:|z-z_c|<R\}$.  Suppose $H$ is holomorphic on a
neighborhood of $\overline D\times\overline\Omega_R$ and
\begin{equation}\label{tr:proper-cluster}
 H(t,z)\ne0\qquad(t\in\overline D,\ |z-z_c|=R).
\end{equation}
The number $n$ of roots in $\Omega_R$, including multiplicity, is then
constant.  We call this a proper finite spatial cluster.

\begin{lemma}[Preparation from contour moments]\label{tr:preparation}
Under \eqref{tr:proper-cluster}, there is a unique monic polynomial
$P(t,z)$ of degree $n$ in $z$, with coefficients holomorphic in $t$,
whose roots are precisely the selected cluster.  Moreover
$H(t,z)=U(t,z)P(t,z)$, where $U$ is holomorphic and nonvanishing on a
neighborhood of the closed product.
\end{lemma}

\begin{proof}
For $1\leq k\leq n$, set
\[
 S_k(t)=\frac1{2\pi i}\oint_{|\zeta-z_c|=R}
           \zeta^k\frac{H_z(t,\zeta)}{H(t,\zeta)}\,d\zeta.
\]
Residues identify these holomorphic functions as the power sums of the
cluster roots.  Newton's identities recursively express the elementary
symmetric functions as polynomials in $S_1,\ldots,S_n$; these are the
coefficients of the required monic $P$.  At each fixed $t$, the quotient
$H/P$ has only removable singularities and is nonzero throughout the
disk.  Its joint holomorphy follows from
\[
 U(t,z)=\frac1{2\pi i}\oint_{|\zeta-z_c|=R}
             \frac{H(t,\zeta)}{P(t,\zeta)(\zeta-z)}\,d\zeta.
\]
For each $t$, this is precisely the removable quotient by Cauchy's
formula.  On a collar of the spatial boundary it agrees with the
ordinary quotient, giving the extension there.  Strict nonvanishing on
the compact wall persists after a small enlargement, which also gives
the asserted neighborhood.  Uniqueness follows from the prescribed
root multiset and monicity.
\end{proof}

Define the finite discriminant by
\begin{equation}\label{tr:discriminant}
 \Delta_P(t)=\prod_{i<j}(z_i(t)-z_j(t))^2,
                   \qquad \Delta_P=1\ \text{if }n\leq1.
\end{equation}
It is a symmetric polynomial in the roots and therefore holomorphic in
the coefficients of $P$, even through collisions and without a global
choice of root labels.

\begin{proposition}[The local heat residue]\label{tr:local-residue}
For the local cluster at an order-$m$ heat event from
Theorem~\ref{tr:splitting}, put $q_m=\binom m2$.  Its discriminant has
\begin{equation}\label{tr:disc-expansion}
 \Delta(t_*+s)=d_m s^{q_m}+O(s^{q_m+1}),\qquad
 d_m=2^{q_m}\prod_{i<j}(\lambda_i-\lambda_j)^2>0.
\end{equation}
Consequently $d\log\Delta$ has residue $q_m$ at the event time.
\end{proposition}

\begin{proof}
Substitute the $m$ root branches \eqref{tr:root-expansion} into
\eqref{tr:discriminant}.  Each squared difference has leading term
$2(\lambda_i-\lambda_j)^2\varepsilon^2$.  Their product is
$d_m\varepsilon^{2q_m}(1+O(\varepsilon))$.  The discriminant is
holomorphic in $s=\varepsilon^2$, so its Taylor series in $\varepsilon$
is even.  This improves the remainder to $O(\varepsilon^2)$ and proves
\eqref{tr:disc-expansion}.  Taking the logarithmic derivative gives the
residue.  The coefficient is positive because the Hermite roots are
real and distinct; the original scalar factor $c_m$ does not enter a
monic-root discriminant.
\end{proof}

\begin{theorem}[Finite complex event charge]\label{tr:complex-charge}
Assume \eqref{tr:proper-cluster}, the heat equation, and
\[
 (H,H_z)\ne(0,0)\qquad
             \text{on }\partial D\times\overline\Omega_R.
\]
Then there are finitely many events in $D\times\Omega_R$, and
\begin{equation}\label{tr:charge-residue}
 Q_{\mathbb C}:=\frac1{2\pi i}\oint_{\partial D}d\log\Delta_P
  =\sum_{(t,z)\in D\times\Omega_R:\,H=H_z=0}\binom{m(t,z)}2.
\end{equation}
The boundary of $D$ has its positive complex orientation.  In particular,
all interior complex event residues have the same positive sign.
\end{theorem}

\begin{proof}
The hypotheses make $\Delta_P$ nonzero on $\partial D$.  Thus it is not
identically zero and has finitely many zeros in the closed time domain.
At a zero time, the spatial cluster is finite.  Small disjoint spatial
disks about its distinct roots give separate monic cluster polynomials.
The discriminant of their product is the product of their discriminants
times the squared pairwise resultants.  The latter are nonzero at the
chosen time because roots in distinct disks do not coincide.  Hence the
order of the full discriminant is the sum of the local orders.
Proposition~\ref{tr:local-residue} identifies each order as
$\binom m2$.  The scalar argument principle now proves
\eqref{tr:charge-residue}.
\end{proof}

\begin{proposition}[Agreement with the complex boundary degree]
\label{tr:degree-agreement}
Under the hypotheses of Theorem~\ref{tr:complex-charge}, give
$D\times\Omega_R$ the orientation
$d\operatorname{Re}t\wedge d\operatorname{Im}t
 \wedge d\operatorname{Re}z\wedge d\operatorname{Im}z$.
Then
\begin{equation}\label{tr:degree-formula}
 \deg\left(\frac{(H,H_z)}{\sqrt{|H|^2+|H_z|^2}}\right)
                  \bigg|_{\partial(D\times\Omega_R)}
       =\frac1{2\pi i}\oint_{\partial D}d\log\Delta_P.
\end{equation}
\end{proposition}

\begin{proof}
If $n=0$, then $H$ has no zero on the whole product, the normalized
jet extends across it, and both sides are zero.  Assume henceforth
$n\geq1$.  Here we use the elementary degree facts that a boundary map into the
unit sphere has homotopy-invariant integer degree, and that its degree
is the signed number of inverse images of a nearby regular value.
The latter follows by excising neighborhoods of the inverse images and
applying Stokes' theorem to the normalized sphere volume form; each
invertible derivative contributes the sign of its real determinant.
These facts apply to the piecewise smooth product boundary by rounding
corners in its nonvanishing collar.

By preparation,
\[
 \binom H{H_z}=
 \begin{pmatrix}U&0\\ U_z&U\end{pmatrix}\binom P{P_z}.
\]
Since the product of two disks or a rectangle and a disk is simply
connected and $U$ is nonzero, $U$ has a holomorphic logarithm there.
The invertible matrices
\[
 M_v=\begin{pmatrix}e^{v\log U}&0\\ vU_z&e^{v\log U}\end{pmatrix},
                         \qquad 0\leq v\leq1,
\]
give a nonzero boundary homotopy from $(P,P_z)$ to $(H,H_z)$.
It suffices to compute the degree of $F=(P,P_z)$.

Choose a regular target value $(\eta,\xi)$ with $\xi\ne0$ and norm
smaller than $\min_{\partial(D\times\Omega_R)}|F|$, and small enough
for the Rouch\'e comparison below.  Such choices exist by the regular-value theorem, and they
have no preimage on the product boundary.  Uniform Rouch\'e comparison
on the spatial wall keeps every root of $P(t,\cdot)-\eta$ inside
$\Omega_R$ for all $t\in\overline D$.  Consider the scalar resultant
\[
 R_{\eta,\xi}(t)=\operatorname{Res}_z(P(t,z)-\eta,P_z(t,z)-\xi).
\]
Its zeros are exactly the time coordinates of the inverse images of
$(\eta,\xi)$.  At such an inverse image $P_z=\xi\ne0$, so
$P(t,z)-\eta=0$ has a local branch $z=r(t)$.  The corresponding factor
$g(t)=P_z(t,r(t))-\xi$ has
\[
 g'(t)=P_{zt}-P_{zz}P_t/\xi
           =-\det(D_{(t,z)}F)/\xi\ne0.
\]
Thus this solution contributes one to the zero order of the resultant;
several solutions at the same time contribute additively.  All other
local factors are nonzero.  It follows that the argument-principle
count for $R_{\eta,\xi}$ equals the number of regular inverse images.
Every one contributes $+1$ to degree because the real Jacobian of a
complex-linear invertible derivative is
$|\det_{\mathbb C}DF|^2>0$.

At target zero the resultant is
\[
 R_{0,0}(t)=\prod_iP_z(t,z_i)
          =(-1)^{n(n-1)/2}\Delta_P(t).
\]
The formula holds first for distinct roots and then everywhere by
holomorphy.  Small target perturbations keep the resultant nonzero on
$\partial D$, so its winding is unchanged.  The number of inverse
images, and hence the degree, equals the winding of $\Delta_P$, as
asserted.
\end{proof}

\begin{remark}[Dimension and multiplicity distinctions]\label{tr:distinct-charges}
The integral \eqref{tr:complex-count} counts ordinary zeros on a spatial
complex curve.  The real angular current \eqref{tr:angular-form} is a
one-form on a two-dimensional real spacetime and counts births with
weight $\lfloor m/2\rfloor$.  Formula \eqref{tr:degree-formula} is a
degree on a three-dimensional boundary in complex two-space and has
weight $\binom m2$.  Its reduction to a one-dimensional time contour is
through the discriminant, not through the ordinary divisor count.
The two event weights agree for $m=2$ but already differ for $m=3$.
\end{remark}

\subsection{Cap decompositions and endpoint terms}

\begin{proposition}[The complete scalar cap balance]\label{tr:complex-caps}
Let $D=(a,b)+i(-Y,Y)$ satisfy Theorem~\ref{tr:complex-charge}.  Define
forward vertical cap contributions by
\[
 C_t=\frac1{2\pi}\int_{-Y}^{Y}
             \frac{d}{ds}\arg\Delta_P(t+is)\,ds,
                        \qquad t=a,b,
\]
using any continuous argument along the indicated nonzero segment.
Let $S$ be the sum of the argument increments, divided by $2\pi$, on
the lower horizontal edge traversed from $a-iY$ to $b-iY$ and the upper
horizontal edge traversed from $b+iY$ to $a+iY$.  Then
\begin{equation}\label{tr:complex-cap-balance}
                      Q_{\mathbb C}=C_b-C_a+S.
\end{equation}
Subtracting $d\log\Delta_0$ for a holomorphic nonvanishing reference
$\Delta_0$ on a neighborhood of $\overline D$ changes the individual
edge contributions but leaves their total unchanged.
\end{proposition}

\begin{proof}
The argument increment around a nonzero closed loop is the imaginary
part of its logarithmic-derivative integral.  The real part is the
increment of $\log|\Delta_P|$ and telescopes to zero.  Splitting the
positively oriented boundary into its four edges gives the stated
signs.  A nonvanishing holomorphic reference on the simply connected
rectangle has a holomorphic logarithm, whose closed-contour integral
is zero.
\end{proof}

Consequently, for a family of complete domains carrying the same
interior charge, genuinely proved limits $C_b\to0$ and $S\to0$ would
imply $C_a\to-Q_{\mathbb C}$.  This implication is a transport identity,
not a proof of either limit.  A late cap with only simple roots merely
has $\Delta_P\ne0$; its argument may still vary.  A reference defined
only locally, or one with its own zeros or poles, instead brings its
transition or divisor terms and is not covered by the final assertion
of Proposition~\ref{tr:complex-caps}.

\begin{proposition}[An explicitly retained endpoint indentation]
\label{tr:indentation}
Suppose a scalar discriminant has an event at a boundary time $t_0$,
with $\Delta(t)=(t-t_0)^q g(t)$, $q\geq1$, and $g$ holomorphic and
nonzero near $t_0$.  On an oriented circular arc
$t=t_0+\varepsilon e^{i\theta}$, with $\theta$ running from
$\theta_0$ to $\theta_1$ in the declared orientation,
\begin{equation}\label{tr:arc-term}
 \frac1{2\pi i}\int_{\mathrm{arc}}d\log\Delta
        =\frac{q(\theta_1-\theta_0)}{2\pi}+O(\varepsilon).
\end{equation}
A clockwise semicircular indentation therefore contributes
$-q/2+O(\varepsilon)$.  The argument principle applies to the indented
domain with this arc included, counting only the events left in its
interior.
\end{proposition}

\begin{proof}
Write $d\log\Delta=q\,dt/(t-t_0)+(g'/g)dt$.  The first integral is
$iq(\theta_1-\theta_0)$; the second is $O(\varepsilon)$ because
$g'/g$ is bounded and the arc has length $O(\varepsilon)$.  Division
by $2\pi i$ gives \eqref{tr:arc-term}.  The final assertion is the
ordinary argument principle on the boundary-regular indented domain.
\end{proof}

\begin{remark}[No silent endpoint convention]\label{tr:endpoint-warning}
A real multiple zero on the initial cap violates the boundary-jet
hypothesis of Theorem~\ref{tr:real-balance}.  One may start at an earlier
regular time, or excise a prescribed small neighborhood of the endpoint
and retain the resulting oriented arc in the real angular integral.
Stokes' theorem remains exact on this excised domain.  Its real angular
arc is not automatically given by the scalar discriminant formula
\eqref{tr:arc-term}; these are different currents on different spaces.
In particular, no universal ``half charge'' for an unspecified endpoint
regularization is assumed here.

An order-$m$ real zero on the initial cap, counted with its spatial
multiplicity, already contributes $m$.  Its $m$ simple real descendants
at positive time also contribute $m$.  Thus ordinary multiplicity
counts at these two times do not detect the initial multiplicity.
Moreover \eqref{tr:single-cap} cannot be evaluated at that endpoint by
replacing its regular-root count with a multiplicity count: the angular
current itself is undefined at the joint zero.  This is the exact
reason that an all-zero cap difference alone is not an endpoint
simplicity theorem.
\end{remark}

\subsection{Exact controls}

For $H(t,z)=z^2-2t$, a sufficiently wide proper spatial window has two
complex zeros on either regular time cap, whereas its real counts are
zero before $t=0$ and two after it.  Its real charge is one and
$\Delta=8t$ gives complex charge one.  For
$H(t,z)=z^3-6tz$, the real counts are one and three, so its real charge
is again one, while $\Delta=864t^3$ gives complex charge three.
Both controls satisfy the analytic and heat hypotheses exactly.
They verify the normalization and illustrate that none of the
conservation identities excludes an event.  No numerical approximation
is used in any proof of this section.

\begin{remark}[Scope]\label{tr:scope}
These identities hold for every heat family satisfying the stated
finite-domain hypotheses.  The explicit polynomial controls show that
they are compatible with multiple-zero events.  An exclusion theorem
requires an additional analytic estimate on the complete boundary
integral, including every cap, wall, and endpoint term, that contradicts
the positive event sum.  Neither simplicity of a late cap nor
conservation of its ordinary complex multiplicity supplies that estimate.
\end{remark}
